\documentclass[11pt]{article}
\usepackage[margin=1in]{geometry}
\usepackage{amsmath,amssymb,amsthm}
\usepackage{enumitem}
\usepackage{microtype}
\usepackage[T1]{fontenc}
\usepackage[utf8]{inputenc}

\setlength{\parindent}{0pt}
\setlength{\parskip}{0.65em}
\raggedbottom
\setlist[enumerate]{leftmargin=1.7em,itemsep=0.35em,topsep=0.35em}

\newcommand{\norm}[1]{\left\lVert #1\right\rVert}
\newcommand{\dd}{\,d}

\theoremstyle{definition}
\newtheorem*{proposition}{Proposition}

\title{Decay of Heat-Scale Residual Refill at High Frequency}
\author{}
\date{}

\begin{document}
\maketitle

\section*{The quantity being measured}

The incompressible three-dimensional Navier-Stokes equation is
\[
\partial_t u+(u\cdot\nabla)u+\nabla p-\nu\Delta u=0,
\qquad
\nabla\cdot u=0,
\]
where \(u\) is velocity, \(p\) is pressure, and \(\nu>0\) is viscosity. A
finite-time concentration scenario means that, as time approaches a possible
terminal time \(T_*\), the solution's active part is forced into smaller and
smaller spatial scales near the terminal region.

Annular or funnel windows are used to watch that terminal region without
freezing the observation at one fixed scale. Each window has a terminal end
time and a short backward time length. As the windows approach \(T_*\), they
look for activity that survives closer and closer to the terminal point.

After a Littlewood-Paley decomposition, smaller spatial scales appear as larger
dyadic frequencies. The quantity defined below measures whether residual source
can keep refilling those high-frequency shells once heat damping has acted.

Let \(\Delta_j\) be the Littlewood-Paley projection onto the shell whose
frequency size is about \(2^j\), and write
\[
u_j=\Delta_j u.
\]
Large \(j\) means high frequency and small spatial scale.

Viscosity damps shell \(j\) at rate comparable to
\[
\nu2^{2j},
\]
where \(\nu>0\) is the viscosity. The natural heat time for shell \(j\) is
then of order
\[
2^{-2j}.
\]
A source feeding shell \(j\) matters at high frequency through the part
that survives this heat-scale damping.

Let \(\Phi_j^+(\sigma)\) be the nonnegative residual source feeding shell \(j\)
at time \(\sigma\). For two times \(t_0<t<T_*\), define the heat-scale residual
refill by
\[
D_j^{\rm res}(t;t_0)
=
2^{2j}\int_{t_0}^{t}
e^{-c_0\nu2^{2j}(t-\sigma)}
\Phi_j^+(\sigma)\dd\sigma .
\]
Here \(c_0>0\) is a fixed heat-kernel constant. The exponential factor records
heat damping from time \(\sigma\) to time \(t\). The outside factor \(2^{2j}\)
normalizes the integral to the heat time scale of the shell. Thus
\(D_j^{\rm res}(t;t_0)\) measures how much residual source still refills shell
\(j\) after heat damping.

\section*{Dyadic input}

The following dyadic source estimate is the input to the decay argument. Let
\(s>5/2\), and assume that the
velocity field has the uniform Sobolev bound
\[
\sup_{t<T_*}\norm{u(t)}_{H^s}\le M,
\]
where \(M\) is fixed. Assume also that the residual source obeys
\[
\Phi_j^+(t)
\le
C_s\,2^{2j}
\norm{u_j(t)}_2
\norm{\Delta_j\mathbb P\nabla\cdot(u\otimes u)(t)}_2 .
\]
Here \(C_s\) depends on \(s\), and \(\mathbb P\) is the Leray projection onto
divergence-free vector fields. All remaining estimates below come from this
source bound and the \(H^s\) hypothesis.

\begin{proposition}[Decay of residual refill under an \(H^s\) bound]
Under the assumptions above, there is a constant \(C_s'>0\), independent of
\(j\), such that
\[
D_j^{\rm res}(t;t_0)
\le
C_s' M^3 2^{(3-2s)j}.
\]
In particular,
\[
D_j^{\rm res}(t;t_0)\to0
\quad\text{as }j\to\infty.
\]
\end{proposition}

\section*{Proof}

The proof is a power count. Each step records the power of \(2^j\) that remains
at shell \(j\).

First, the \(H^s\) bound controls one shell. On shell \(j\), applying \(s\)
derivatives costs \(2^{sj}\). Hence
\[
2^{sj}\norm{u_j(t)}_2\le M.
\]
Dividing by \(2^{sj}\) gives
\[
\norm{u_j(t)}_2\le M2^{-sj}.
\]

Second, the nonlinear term has one derivative. Since \(s>5/2\), the space
\(H^s(\mathbb R^3)\) is an algebra. Thus
\[
\norm{u\otimes u}_{H^s}
\le
C_s\norm{u}_{H^s}^2
\le
C_sM^2.
\]
Taking the \(j\)-th shell gives
\[
\norm{\Delta_j(u\otimes u)}_2
\le
C_sM^2 2^{-sj}.
\]
The operator \(\nabla\cdot\) contributes one derivative, and one derivative on
shell \(j\) costs \(2^j\). The Leray projection \(\mathbb P\) has order zero.
Thus
\[
\norm{\Delta_j\mathbb P\nabla\cdot(u\otimes u)(t)}_2
\le
C_sM^2 2^j2^{-sj}.
\]
Equivalently,
\[
\norm{\Delta_j\mathbb P\nabla\cdot(u\otimes u)(t)}_2
\le
C_sM^2 2^{(1-s)j}.
\]

Insert the two shell bounds into the source estimate:
\[
\Phi_j^+(t)
\le
C_s\,2^{2j}
\bigl(M2^{-sj}\bigr)
\bigl(M^2 2^{(1-s)j}\bigr).
\]
The constants multiply as
\[
M\cdot M^2=M^3.
\]
The dyadic powers multiply as
\[
2^{2j}2^{-sj}2^{(1-s)j}
=
2^{(3-2s)j}.
\]
Hence
\[
\Phi_j^+(t)\le C_sM^3 2^{(3-2s)j}.
\]

Now return to the definition of \(D_j^{\rm res}\):
\[
D_j^{\rm res}(t;t_0)
=
2^{2j}\int_{t_0}^{t}
e^{-c_0\nu2^{2j}(t-\sigma)}
\Phi_j^+(\sigma)\dd\sigma.
\]
Using the source bound just proved,
\[
D_j^{\rm res}(t;t_0)
\le
2^{2j}C_sM^3 2^{(3-2s)j}
\int_{t_0}^{t}
e^{-c_0\nu2^{2j}(t-\sigma)}\dd\sigma.
\]

The remaining integral is controlled by the full heat tail:
\[
\int_{t_0}^{t}
e^{-c_0\nu2^{2j}(t-\sigma)}\dd\sigma
\le
\int_0^\infty e^{-c_0\nu2^{2j}\rho}\dd\rho.
\]
The last integral is explicit:
\[
\int_0^\infty e^{-c_0\nu2^{2j}\rho}\dd\rho
=
\frac{1}{c_0\nu2^{2j}}.
\]
Substituting this gives
\[
D_j^{\rm res}(t;t_0)
\le
2^{2j}C_sM^3 2^{(3-2s)j}
\frac{1}{c_0\nu2^{2j}}.
\]
The two factors \(2^{2j}\) cancel:
\[
D_j^{\rm res}(t;t_0)
\le
\frac{C_s}{c_0\nu}M^3 2^{(3-2s)j}.
\]
Set
\[
C_s'=\frac{C_s}{c_0\nu}.
\]
This constant is independent of \(j\), so
\[
D_j^{\rm res}(t;t_0)
\le
C_s'M^3 2^{(3-2s)j}.
\]

It remains to check the exponent. Since \(s>5/2\),
\[
2s>5.
\]
Subtracting from \(3\) gives
\[
3-2s<3-5=-2.
\]
In particular,
\[
3-2s<0.
\]
Set
\[
a=2s-3.
\]
Then \(a>0\) and \(3-2s=-a\). Hence
\[
2^{(3-2s)j}=2^{-aj}=\frac{1}{2^{aj}}.
\]
As \(j\to\infty\), \(2^{aj}\to\infty\), so
\[
2^{(3-2s)j}\to0.
\]
The constants \(C_s'\) and \(M\) are fixed while \(j\to\infty\). Thus
\[
C_s'M^3 2^{(3-2s)j}\to0.
\]
The bound
\[
0\le D_j^{\rm res}(t;t_0)\le C_s'M^3 2^{(3-2s)j}
\]
then gives
\[
D_j^{\rm res}(t;t_0)\to0.
\]
This proves the proposition.

\section*{Contradiction with a selected terminal witness}

Suppose an annular or funnel selection produces time windows and shell indices
\[
j_m\to\infty.
\]
The index \(m\) labels the selected window. Its end time is \(t_m\), its length
is \(\tau_m\), and its start time is \(t_m-\tau_m\). The selected residual
refill is
\[
D_{j_m}^{\rm res}(t_m;t_m-\tau_m).
\]
The proposition gives
\[
D_{j_m}^{\rm res}(t_m;t_m-\tau_m)
\le
C_s'M^3 2^{(3-2s)j_m}.
\]
Since \(j_m\to\infty\) and \(3-2s<0\), the right hand side tends to zero. For
each fixed \(\eta>0\), there is an index \(m_0\) such that, for every
\(m\ge m_0\),
\[
C_s'M^3 2^{(3-2s)j_m}<\eta.
\]
For those \(m\),
\[
D_{j_m}^{\rm res}(t_m;t_m-\tau_m)<\eta.
\]

A fixed positive witness asserts the opposite lower bound along the same
selected sequence:
\[
D_{j_m}^{\rm res}(t_m;t_m-\tau_m)\ge\eta.
\]
For large \(m\), the same number is smaller than \(\eta\) and at least
\(\eta\). The contradiction shows that a uniform \(H^s\) bound with
\(s>5/2\) cannot coexist with a fixed positive residual refill witness at
frequencies \(j_m\to\infty\).

\end{document}
